\section*{الفصل \ref{ch:b1:complex} --- الأعداد العقدية}

\begin{solution}{exo:b1:complex:1}
$1 + \iu = \sqrt 2\, \eu^{\iu\pi/4}$؛
$\;\sqrt 3 - \iu = 2\, \eu^{-\iu\pi/6}$؛
$\;-5 = 5\, \eu^{\iu\pi}$؛
\[
\frac{1 + \iu}{\sqrt 3 - \iu}
= \frac{\sqrt 2}{2}\, \eu^{\iu(\pi/4 + \pi/6)}
= \frac{\sqrt 2}{2}\, \eu^{5\iu\pi/12}.
\]
وبحساب القسمة نفسها جبريًا (بالضرب في \omterm{def:b1:complex:field}{المرافق}):
\[
\frac{(1 + \iu)(\sqrt 3 + \iu)}{4}
= \frac{(\sqrt 3 - 1) + \iu(\sqrt 3 + 1)}{4}.
\]
وبالمطابقة مع $\frac{\sqrt 2}{2}(\cos\frac{5\pi}{12} +
\iu\sin\frac{5\pi}{12})$:
\[
\cos\frac{5\pi}{12} = \frac{\sqrt 6 - \sqrt 2}{4},
\qquad
\sin\frac{5\pi}{12} = \frac{\sqrt 6 + \sqrt 2}{4}.
\]
\end{solution}

\begin{solution}{exo:b1:complex:2}
لدينا $(1+\iu)^2 = 2\iu$، ومنه $(1+\iu)^{20} = (2\iu)^{10} = 2^{10}\,\iu^{10}
= 1024 \times (-1) = -1024$.

وفي الشكل الأسّي $(1+\iu)^n = 2^{n/2}\, \eu^{\iu n\pi/4}$، وهو حقيقي إذا
وفقط إذا كان $\sin\frac{n\pi}{4} = 0$، أي $4 \mid n$. إذن $(1+\iu)^n
\in \R$ تحديدًا من أجل مضاعفات $4$ (والقيمة $(-4)^{n/4}$).
\end{solution}

\begin{solution}{exo:b1:complex:3}
$\abs{z - 2} = \abs{z + \iu}$: تساوي البعد عن النقطتين $2$ و
$-\iu$ --- أي محور القطعة الواصلة بين $(2, 0)$
و $(0, -1)$.

$\abs{z - 1} = 2$: الدائرة التي مركزها $1$ ونصف قطرها $2$.

$\frac{z-1}{z+1} \in \iu\R$: حسب
\cref{prop:b1:complex:geometry}~(3)، تحقق النقاط $M(z)$ و $A(1)$ و
$B(-1)$ ما يلي: المستقيمان $MA$ و $MB$ متعامدان (أو $z = 1$، حيث
تكون القسمة $0 \in \iu\R$). \omterm{def:b1:logic:sets}{والمجموعة} هي الدائرة ذات القطر
$[-1, 1]$ (أي دائرة الوحدة)، منقوصةً منها النقطة $-1$ حيث تكون
القسمة غير معرَّفة. \emph{وللتحقق بالحساب:} $z = \eu^{\iu\theta}$ يعطي
$\frac{z-1}{z+1} = \frac{\eu^{\iu\theta/2}(\eu^{\iu\theta/2} -
\eu^{-\iu\theta/2})}{\eu^{\iu\theta/2}(\eu^{\iu\theta/2} +
\eu^{-\iu\theta/2})} = \iu\tan\frac\theta2 \in \iu\R$.
\end{solution}

\begin{solution}{exo:b1:complex:4}
التخطيط:
\[
\sin^4\theta
= \Bigl(\frac{\eu^{\iu\theta} - \eu^{-\iu\theta}}{2\iu}\Bigr)^{\!4}
= \frac{\eu^{4\iu\theta} - 4\eu^{2\iu\theta} + 6 - 4\eu^{-2\iu\theta}
+ \eu^{-4\iu\theta}}{16}
= \frac{\cos 4\theta - 4\cos 2\theta + 3}{8}.
\]
والنشر: بصيغة دي موافر، $\cos 4\theta = \Re\bigl((c +
\iu s)^4\bigr) = c^4 - 6c^2 s^2 + s^4$ حيث $c = \cos\theta$ و$s =
\sin\theta$؛ وبتعويض $s^2 = 1 - c^2$:
\[
\cos 4\theta = c^4 - 6c^2(1 - c^2) + (1 - c^2)^2
= 8c^4 - 8c^2 + 1 .
\]
\end{solution}

\begin{solution}{exo:b1:complex:5}
$8\iu = 8\,\eu^{\iu\pi/2}$، ومنه فالجذور التكعيبية هي $2\,\eu^{\iu(\pi/6 +
2k\pi/3)}$ حيث $k = 0, 1, 2$:
\[
z_0 = 2\eu^{\iu\pi/6} = \sqrt 3 + \iu,\quad
z_1 = 2\eu^{5\iu\pi/6} = -\sqrt 3 + \iu,\quad
z_2 = 2\eu^{3\iu\pi/2} = -2\iu :
\]
وهي مثلث متساوي الأضلاع على الدائرة ذات نصف القطر $2$.

ولدينا $-4 = 4\eu^{\iu\pi}$، ومنه فحلول $z^4 = -4$ هي $\sqrt 2\,
\eu^{\iu(\pi/4 + k\pi/2)}$: أي $\;1 + \iu$ و $-1 + \iu$ و $-1 - \iu$
و $1 - \iu$. وبازدواج الجذرين المترافقين:
\[
X^4 + 4 = \bigl(X^2 - 2X + 2\bigr)\bigl(X^2 + 2X + 2\bigr),
\]
لأن $(X - (1+\iu))(X - (1-\iu)) = X^2 - 2X + 2$ وكذلك من أجل
الزوج الآخر. (انشر للتحقق.)
\end{solution}

\begin{solution}{exo:b1:complex:6}
$C_n + \iu S_n = \sum_{k=0}^{n} \eu^{\iu k\theta} =
\frac{\eu^{\iu(n+1)\theta} - 1}{\eu^{\iu\theta} - 1}$ (مجموع هندسي
أساسه $\eu^{\iu\theta} \neq 1$). وبتعميل نصف الزاوية
(\cref{met:b1:complex:trig}~(4)) في البسط والمقام:
\[
\frac{2\iu\sin\frac{(n+1)\theta}{2}\, \eu^{\iu(n+1)\theta/2}}
{2\iu\sin\frac{\theta}{2}\, \eu^{\iu\theta/2}}
= \frac{\sin\frac{(n+1)\theta}{2}}{\sin\frac{\theta}{2}}\,
\eu^{\iu n\theta/2}.
\]
وبأخذ الجزأين الحقيقي والتخيلي:
\[
C_n = \frac{\sin\frac{(n+1)\theta}{2}}{\sin\frac{\theta}{2}}
\cos\frac{n\theta}{2},
\qquad
S_n = \frac{\sin\frac{(n+1)\theta}{2}}{\sin\frac{\theta}{2}}
\sin\frac{n\theta}{2}.
\]
\end{solution}

\begin{solution}{exo:b1:complex:7}
المميِّز: $\Delta = (3 + 4\iu)^2 - 4(-1 + 5\iu) = 9 + 24\iu - 16 +
4 - 20\iu = -3 + 4\iu$. والجذران التربيعيان للعدد $-3 + 4\iu$: حُلَّ $x^2 - y^2
= -3$ و $2xy = 4$ و $x^2 + y^2 = 5$؛ عندئذ $x^2 = 1$ و $y = 2/x$، وبالإشارتين
نفسيهما: $\delta = \pm(1 + 2\iu)$. ومنه
\[
z = \frac{(3 + 4\iu) \pm (1 + 2\iu)}{2}
\in \{\, 2 + 3\iu,\; 1 + \iu \,\}.
\]
\emph{وللتحقق:} المجموع $= 3 + 4\iu$ والجداء $(2+3\iu)(1+\iu) = -1 +
5\iu$، كما تقتضيه المعاملات.
\end{solution}

\begin{solution}{exo:b1:complex:8}
\begin{enumerate}
  \item \cref{prop:b1:complex:sumroots} مع $n = 5$.
  \item $u + v = \omega + \omega^2 + \omega^3 + \omega^4 = -1$ حسب (1).
        وأمّا الجداء فانشر وأرجِع الأسس بترديد $5$:
        \[
        uv = (\omega + \omega^4)(\omega^2 + \omega^3)
        = \omega^3 + \omega^4 + \omega^6 + \omega^7
        = \omega^3 + \omega^4 + \omega + \omega^2 = -1 .
        \]
        إذن مجموع $u$ و $v$ هو $-1$ وجداؤهما $-1$: فهما
        جذرا كثير الحدود $X^2 + X - 1$.
  \item لدينا $u = \omega + \conj\omega = 2\cos\frac{2\pi}{5} > 0$ (فالزاوية
        حادّة)، والجذر الموجب لكثير الحدود $X^2 + X - 1$ هو
        $\frac{-1 + \sqrt 5}{2}$. ومنه $\cos\frac{2\pi}{5} =
        \frac{\sqrt 5 - 1}{4}$.
\end{enumerate}
\end{solution}

\begin{solution}{exo:b1:complex:9}
انشر المربعين كما في برهان
\cref{prop:b1:complex:rules}~(4):
\[
\abs{z + w}^2 = \abs z^2 + \abs w^2 + 2\Re(z\conj w),
\qquad
\abs{z - w}^2 = \abs z^2 + \abs w^2 - 2\Re(z\conj w),
\]
ثم اجمع. وهندسيًا، $\abs{z+w}$ و $\abs{z-w}$ هما طولا
قطرَي متوازي الأضلاع ذي الرؤوس $0, z, z+w, w$،
بينما $\abs z$ و $\abs w$ هما طولا الضلعين: فمجموع
مربعَي القطرين يساوي مجموع مربعات الأضلاع الأربعة.
\end{solution}

\begin{solution}{exo:b1:complex:10}
ليكن $j = \eu^{2\iu\pi/3}$، ومنه $j^2 + j + 1 = 0$ و$\eu^{\iu\pi/3} =
-j^2$ و $\eu^{-\iu\pi/3} = -j$. ويكون المثلث متساوي الأضلاع مباشرًا إذا وفقط إذا كان
$c - a = -j^2 (b - a)$، وغير مباشر إذا وفقط إذا كان $c - a = -j(b - a)$
(\cref{ex:b1:complex:equilateral}).

تأمّل $P = a + jb + j^2 c$ و $Q = a + j^2 b + jc$. وباستعمال $1 +
j^2 = -j$ و $j^3 = 1$:
\[
c - a + j^2(b - a) = c + j^2 b - (1 + j^2)\,a = c + j^2 b + ja
= j\,(a + jb + j^2 c) = jP,
\]
فتُقرأ الحالة المباشرة $jP = 0 \iff P = 0$؛ والحساب نفسه
مع $j$ بدل $j^2$ يعطي $c - a + j(b - a) = j^2 Q$، فتُقرأ
الحالة غير المباشرة $Q = 0$. ومنه: يكون المثلث متساوي الأضلاع (بأيّ توجيه)
$\iff PQ = 0$. وبالنشر، مع $j + j^2 = -1$:
\[
PQ = a^2 + b^2 + c^2 + (j + j^2)(ab + bc + ca)
= a^2 + b^2 + c^2 - (ab + bc + ca).
\]
إذن يكون المثلث متساوي الأضلاع إذا وفقط إذا كان $a^2 + b^2 + c^2 = ab +
bc + ca$.
\end{solution}

\begin{solution}{exo:b1:complex:11}
بما أن الأعداد $\omega^k$ حيث $k = 0, \dots, n-1$ هي بالضبط الجذور $n$
لكثير الحدود $X^n - 1$ (\cref{thm:b1:complex:roots})، وأن كثير الحدود
واحديّ:
\[
X^n - 1 = \prod_{k=0}^{n-1} (X - \omega^k)
= (X - 1) \prod_{k=1}^{n-1} (X - \omega^k).
\]
وبالقسمة على $X - 1$: $\;1 + X + \dots + X^{n-1} = \prod_{k=1}^{n-1}
(X - \omega^k)$. وبتعويض $X = 1$ نجد $P = n$.

والآن $1 - \omega^k = -\eu^{\iu k\pi/n}\bigl(\eu^{\iu k\pi/n} -
\eu^{-\iu k\pi/n}\bigr) = -2\iu\,\eu^{\iu k\pi/n} \sin\frac{k\pi}{n}$،
فبأخذ المقاييس في $P = n$ (وكل $\sin\frac{k\pi}n > 0$ من أجل $1 \leq k
\leq n-1$):
\[
n = \abs P = \prod_{k=1}^{n-1} 2\sin\frac{k\pi}{n}
= 2^{n-1} \prod_{k=1}^{n-1} \sin\frac{k\pi}{n},
\qquad\text{ومنه}\qquad
\prod_{k=1}^{n-1} \sin\frac{k\pi}{n} = \frac{n}{2^{n-1}} .
\]
\end{solution}

\begin{solution}{exo:b1:complex:12}
العدد $z = 1$ ليس حلًّا (إذ $2^n \neq 0$)، فتكافئ المعادلة
$\bigl(\frac{z+1}{z-1}\bigr)^n = 1$، أي
$\frac{z+1}{z-1} = \omega^k$ حيث $\omega = \eu^{2\iu\pi/n}$ و
$k \in \{0, \dots, n-1\}$. والقيمة $k = 0$ مستبعدة (إذ $z + 1 =
z - 1$ مستحيل). ومن أجل $1 \leq k \leq n - 1$، يعطي حلّ
$z + 1 = \omega^k(z - 1)$ أن $z(1 - \omega^k) = -1 - \omega^k$،
ومنه، مع $\varphi = \frac{2k\pi}{n}$ وبتعميلَي نصف
الزاوية $1 + \eu^{\iu\varphi} = 2\cos\frac\varphi2\,
\eu^{\iu\varphi/2}$ و$\eu^{\iu\varphi} - 1 = 2\iu\sin\frac\varphi2
\,\eu^{\iu\varphi/2}$:
\[
z_k = \frac{1 + \omega^k}{\omega^k - 1}
= \frac{2\cos\frac{k\pi}{n}}{2\iu\,\sin\frac{k\pi}{n}}
= -\iu\,\frac{\cos\frac{k\pi}{n}}{\sin\frac{k\pi}{n}} ,
\]
وهي تخيلية بحتة كما ادُّعي. \omterm{def:b1:logic:map}{والتطبيق} $t \mapsto \cos t/\sin t$
متباين على $\intoo0\pi$ (فهو متناقص تمامًا)، فتكون القيم
$z_k$ التي عددها $n - 1$ متمايزة مثنى مثنى: أي أن للمعادلة، وهي من الدرجة
$n - 1$ بعد النشر (إذ تختصر حدود $z^n$)، هذه الحلول
$n - 1$ بالضبط.
\end{solution}

\begin{solution}{pb:b1:complex:1}
\textbf{1.} $j = \cos\frac{2\pi}3 + \iu\sin\frac{2\pi}3 = -\frac12
+ \iu\frac{\sqrt3}2$؛ و$j^2 = \eu^{4\iu\pi/3} = -\frac12 -
\iu\frac{\sqrt3}2 = \conj j$؛ و $j^3 = 1$؛ و $1 + j + j^2 = 0$ (وهو مجموع
الجذور التكعيبية للوحدة، \cref{prop:b1:complex:sumroots})؛ و$j^{-1}
= j^2$ (لأن $j \cdot j^2 = 1$). وعلى دائرة الوحدة، تكون $1$ و $j$
و $j^2$ رؤوسًا لمثلث مباشر متساوي الأضلاع.

\textbf{2.} الدوران الذي مركزه $a$ وزاويته $\theta$ يصمد عند $a$
ويدير كل متجهة صادرة عن $a$ بالزاوية $\theta$: $r(z) - a =
\eu^{\iu\theta}(z - a)$، أي $r(z) = a + \eu^{\iu\theta}(z - a)$.
وبتركيب $r(z) = a + \eu^{\iu\theta}(z-a)$ مع $r'(z) = a' +
\eu^{\iu\theta'}(z-a')$:
\[
r' \circ r\,(z) = \eu^{\iu(\theta + \theta')} z +
\text{ثابت},
\]
وهو \omterm{def:b1:logic:map}{تطبيق} على الصورة $Az + B$ حيث $A = \eu^{\iu(\theta+\theta')}$ ذو
\omterm{def:b1:complex:field}{مقياس} $1$. وحسب \cref{prop:b1:complex:similitude}، يكون دورانًا
زاويته $\arg A = \theta + \theta'$ إذا كان $A \neq 1$، ويكون
انسحابًا إذا كان $A = 1$، أي إذا كان $\theta + \theta' \in 2\pi\Z$.

\textbf{3.} تكون $(b, c, p)$ متساوية الأضلاع إذا وفقط إذا كان $\abs{p - b} = \abs{c -
b} = \abs{p - c}$. وبكتابة $q = \frac{p - b}{c - b}$، يقول التساوي
الأول إن $\abs q = 1$ ويقول الثاني إن $\abs{q - 1} = 1$؛
ومعًا يعطيان $q = \eu^{\pm\iu\pi/3}$ (فنقطتا تقاطع
الدائرتين $\abs q = 1$ و $\abs{q - 1} = 1$ هما $\frac12 \pm
\iu\frac{\sqrt3}2$). ومنه $p = b + \eu^{\pm\iu\pi/3}(c - b)$. ومن أجل
$a = 0$ و $b = 1$ و $c = \iu$ (وهو مثلث مباشر): يعطي الاختيار
$\eu^{-\iu\pi/3}$
\[
p = 1 + \Bigl(\tfrac12 - \iu\tfrac{\sqrt3}2\Bigr)(\iu - 1)
= \tfrac{1 + \sqrt3}2\,(1 + \iu) \approx 1.37 + 1.37\,\iu ,
\]
وهي تقع على الجهة البعيدة من المستقيم $BC$ (أي $x + y = 1$) عن $a
= 0$: أي نحو الخارج. ويعطي الاختيار $\eu^{+\iu\pi/3}$ القيمة $p \approx -0.37
- 0.37\,\iu$، وهي على الجهة نفسها التي فيها $a$: أي نحو الداخل.

\textbf{4.} اضرب $P = a + jb + j^2c$ في $j^2$: $j^2 P = j^2 a +
j^3 b + j^4 c = b + jc + j^2 a$ (باستعمال $j^3 = 1$ و $j^4 = j$)؛ واضربه
في $j$: $jP = ja + j^2 b + c$. وهاتان هما المتطابقتان. وإذا كان $P
= 0$ فإن $j^2 P = jP = 0$: أي أن المعيار يتحقق من أجل $(b, c, a)$ و
$(c, a, b)$ كذلك --- وهو عدم التغير بالتبديل الدائري (وهذا لازم:
فالمثلث المتساوي الأضلاع لا يعنيه أيّ رأس يُذكر أولًا).

\textbf{5.} إذا كان $a = b$: $0 = a(1 + j) + j^2 c = -j^2 a + j^2 c$
(باستعمال $1 + j = -j^2$)، ومنه $c = a$. وإذا كان $b = c$: $0 = a + b(j + j^2)
= a - b$، ومنه $a = b$. وإذا كان $a = c$: $0 = a(1 + j^2) + jb = -ja + jb$،
ومنه $a = b$. ففي كل حالة تنطبق النقاط الثلاث.

\textbf{6.} $(a'z + b') \circ (az + b) = a'a\,z + (a'b + b')$ حيث
$a'a \neq 0$: أي الصورة نفسها. ومقلوب $z \mapsto az + b$ هو $z
\mapsto \frac1a z - \frac ba$، وهو على الصورة نفسها كذلك. ومع
\omterm{def:b1:logic:map}{التطبيق} المطابق عنصرًا محايدًا، تكوّن التشابهات المباشرة
زمرة بالنسبة إلى التركيب.

\textbf{7.} يحقق $f(z) = az + b$ العلاقتين $f(z_1) = w_1$ و$f(z_2) =
w_2$ إذا وفقط إذا كان $a z_1 + b = w_1$ و $a z_2 + b = w_2$؛ وبالطرح،
$a(z_2 - z_1) = w_2 - w_1$، ومنه يُفرض أن
\[
a = \frac{w_2 - w_1}{z_2 - z_1} \;(\neq 0), \qquad
b = w_1 - a z_1
\]
وبالعكس فإن هذا الاختيار يفي بالغرض: أي الوجود
والوحدانية.

\textbf{8.} $\abs{f(z) - f(w)} = \abs{a(z - w)} = \abs a\,\abs{z -
w}$: أي أن جميع المسافات تُضرب في $\abs a$. وأمّا الشكل:
\[
\frac{f(c') - f(a')}{f(b') - f(a')} = \frac{a(c' -
a')}{a(b' - a')} = \frac{c' - a'}{b' - a'} .
\]
وإذا تساوى شكلا مثلثين $(a', b', c')$ و $(a'', b'', c'')$،
فليكن $f$ التشابهَ المباشر الوحيد الذي يحقق $f(a') = a''$
و $f(b') = b''$ (السؤال 7)؛ عندئذ يساوي شكل $(a'', b'',
f(c'))$ شكل $(a', b', c')$، ومنه شكل $(a'', b'',
c'')$، والشكل يحدّد الرأس الثالث انطلاقًا من الرأسين
الأولين: أي $f(c') = c''$. وبالعكس ينتج تساوي الشكلين من
عدم التغير المعروض.

\textbf{9.} لدينا $b = f(0) = 1$؛ و $a + b = f(1) = \iu$ يعطي $a = \iu -
1$. والنسبة $\abs a = \sqrt2$، والزاوية $\arg(\iu - 1) =
\frac{3\pi}4$. والنقطة الصامدة:
\[
\zeta = \frac{b}{1 - a} = \frac1{2 - \iu} = \frac{2 + \iu}5 .
\]
إذن $f$ هو التشابه المباشر الذي مركزه $\frac{2+\iu}5$ ونسبته
$\sqrt2$ وزاويته $\frac{3\pi}4$.

\textbf{10.} مع $\alpha + \beta = 1$ و $f(z) = az + b$:
\[
f(\alpha a' + \beta b') = a\alpha a' + a\beta b' + b
= \alpha(a a' + b) + \beta(a b' + b)
= \alpha f(a') + \beta f(b') ,
\]
والمفتاح هو $b = (\alpha + \beta)b$. ومنه فإن الأوساط ($\alpha = \beta =
\frac12$) ومراكز الثقل (بالتكرار) ومراكز نابليون $\alpha b
+ \beta c$ أدناه \emph{متساوقة}: أي أن تحويل
المثلث يحوّلها تبعًا له. وبهذا يكفي البرهان على \omterm{def:b1:logic:statement}{عبارة} لا تتغير
بالتشابه من أجل مثلث واحد حسن الوضع للبرهان عليها من أجل الجميع
--- وهي الرخصة المستعملة في السؤال 21.

\textbf{11.} الرأس الخارجي على $[b, c]$ هو $p_a = b +
\eu^{-\iu\pi/3}(c - b)$ (السؤال 3)، ومنه فالمركز هو
\[
n_a = \frac{b + c + p_a}3
= \frac{2b + c + \frac{1 - \iu\sqrt3}2\,(c - b)}3
= \frac{(3 + \iu\sqrt3)\,b + (3 - \iu\sqrt3)\,c}6
= \alpha b + \beta c ,
\]
حيث $\alpha = \frac{3 + \iu\sqrt3}6$ و$\beta = \frac{3 -
\iu\sqrt3}6 = \conj\alpha$. والحساب نفسه على الضلعين
$[c, a]$ و $[a, b]$ يعطي $n_b = \alpha c + \beta a$ و$n_c =
\alpha a + \beta b$ (بتبديل دائري للأدوار).

\textbf{12.} $\alpha + \beta = \frac{3 + \iu\sqrt3 + 3 -
\iu\sqrt3}6 = 1$. ومنه
\[
\frac{n_a + n_b + n_c}3
= \frac{(\alpha + \beta)(a + b + c)}3 = \frac{a + b + c}3 :
\]
أي أن المثلثين يشتركان في مركز ثقلهما.

\textbf{13.} جمّع بحسب $\alpha$ و $\beta$ وطبّق متطابقات السؤال 4
على $P = a + jb + j^2c$:
\[
n_a + j n_b + j^2 n_c
= \alpha\,(b + jc + j^2 a) + \beta\,(c + ja + j^2 b)
= \alpha\,j^2 P + \beta\,j P
= (\alpha j^2 + \beta j)\,P .
\]

\textbf{14.} مع $j = \frac{-1 + \iu\sqrt3}2$:
\[
\alpha j = \frac{(3 + \iu\sqrt3)(-1 + \iu\sqrt3)}{12}
= \frac{-3 + 3\iu\sqrt3 - \iu\sqrt3 - 3}{12}
= \frac{-3 + \iu\sqrt3}6 = -\beta ,
\]
ومنه $\alpha j + \beta = 0$، فيكون $\alpha j^2 + \beta j = j(\alpha j
+ \beta) = 0$، ويعطي السؤال 13 أن $n_a + jn_b + j^2n_c = 0$ من أجل
كل مثلث $(a, b, c)$. وحسب المعيار (السؤالان 4--5)، تكوّن
المراكز مثلثًا مباشرًا متساوي الأضلاع أو نقطة واحدة:
وهي مبرهنة نابليون.

\textbf{15.} الرأس الداخلي هو $b + \eu^{+\iu\pi/3}(c - b)$، و
يبدّل حسابُ السؤال 11 مع $\eu^{+\iu\pi/3} = \frac{1 +
\iu\sqrt3}2$ بين $\alpha$ و $\beta$: $m_a = \beta b + \alpha
c$ و$m_b = \beta c + \alpha a$ و$m_c = \beta a + \alpha b$. وباستعمال
المتطابقتين المماثلتين $b + j^2c + ja = j\,Q$ و$c + j^2a + jb
= j^2 Q$ من أجل $Q = a + j^2b + jc$:
\[
m_a + j^2 m_b + j m_c
= \beta(b + j^2 c + ja) + \alpha(c + j^2 a + jb)
= (\beta j + \alpha j^2)\, Q = j(\beta + \alpha j)\,Q = 0 ,
\]
لأن $\alpha j = -\beta$ (السؤال 14). إذن تحقق المراكز الداخلية
معيار تساوي الأضلاع \emph{غير المباشر}: أي متساوية الأضلاع لكن
بالتوجيه المعاكس (أو نقطة).

\textbf{16.} حسب السؤال 5 مطبَّقًا على الثلاثية $(n_a, n_b, n_c)$
(وهي تحقق المعيار المباشر)، ينطبق مركزان إذا وفقط إذا
انطبقت الثلاثة؛ وتنطبق الثلاثة إذا وفقط إذا تحقق \emph{المعياران} معًا،
أي إذا تحقق زيادةً على ذلك $n_a + j^2 n_b + j n_c = 0$. احسب كما في
السؤال 13، بالمتطابقتين $b + j^2c + ja = jQ$ و$c + j^2a +
jb = j^2Q$:
\[
n_a + j^2 n_b + j n_c
= \alpha\,jQ + \beta\,j^2 Q = j(\alpha + \beta j)\,Q .
\]
ومباشرةً: $\beta j = \frac{(3 -
\iu\sqrt3)(-1 + \iu\sqrt3)}{12} = \frac{-3 + 3\iu\sqrt3 +
\iu\sqrt3 + 3}{12} = \frac{\iu\sqrt3}3$، ومنه $\alpha + \beta j =
\frac{3 + \iu\sqrt3 + 2\iu\sqrt3}6 = \frac{1 + \iu\sqrt3}2 \neq
0$. إذن ينحلّ مثلث نابليون الخارجي إذا وفقط إذا كان $Q = a +
j^2b + jc = 0$، أي إذا وفقط إذا كان $(a, b, c)$ مثلثًا غير مباشر
متساوي الأضلاع --- وفي تلك الحالة تتجه الإنشاءات «الخارجية» كلها
نحو داخل المنطقة المحيطة بالمثلث وتشترك في مركز واحد.

\textbf{17.} مع $a = 0$ و $b = 1$ و $c = \iu$:
\[
n_a = \alpha + \beta\iu = \frac{(3 + \sqrt3)(1 + \iu)}6,
\qquad
n_b = \alpha\iu = \frac{-\sqrt3 + 3\iu}6,
\qquad
n_c = \beta = \frac{3 - \iu\sqrt3}6 .
\]
ومربعات الأضلاع: $n_a - n_b = \frac{(3 + 2\sqrt3) +
\iu\sqrt3}6$ يعطي $\abs{n_a - n_b}^2 = \frac{(3 + 2\sqrt3)^2 +
3}{36} = \frac{24 + 12\sqrt3}{36} = \frac{2 + \sqrt3}3$؛ و$n_b -
n_c = \frac{(3 + \sqrt3)(-1 + \iu)}6$ يعطي $\abs{n_b - n_c}^2 =
\frac{2(3 + \sqrt3)^2}{36} = \frac{24 + 12\sqrt3}{36}$؛ و$n_c - n_a
= \frac{-\sqrt3 - \iu(3 + 2\sqrt3)}6$ يعطي القيمة نفسها من جديد.
فمربعات الأضلاع الثلاثة كلها تساوي $\frac{2 + \sqrt3}3$: أي أنه متساوي الأضلاع،
كما وُعد.

\textbf{18.} انشر:
\[
(a - b)(c - d) + (a - d)(b - c)
= (ac - ad - bc + bd) + (ab - ac - bd + cd)
= ab - ad - bc + cd ,
\]
و$(a - c)(b - d) = ab - ad - bc + cd$: فهما متساويان.

\textbf{19.} خذ المقاييس في السؤال 18 وطبّق المتراجحة
المثلثية (\cref{prop:b1:complex:rules}~(4)):
\[
AC \cdot BD = \abs{(a-b)(c-d) + (a-d)(b-c)}
\leq \abs{a-b}\,\abs{c-d} + \abs{a-d}\,\abs{b-c}
= AB \cdot CD + AD \cdot BC ,
\]
مع التساوي إذا وفقط إذا وقع المجموعان على نصف مستقيم واحد مبدؤه $0$
(وهي حالة التساوي في المتراجحة المثلثية).

\textbf{20.} بتعميل نصف الزاوية
(\cref{met:b1:complex:trig}~(4)):
$\eu^{\iu\theta} - \eu^{\iu\varphi} = 2\iu\,\sin\frac{\theta -
\varphi}2\;\eu^{\iu(\theta + \varphi)/2}$، وأخذ المقاييس يُلغي
العوامل الأحادية \omterm{def:b1:complex:field}{المقياس}: $\abs{\eu^{\iu\theta} - \eu^{\iu\varphi}}
= 2\,\abs{\sin\frac{\theta - \varphi}2}$.

\textbf{21.} مع $p = \eu^{\iu\theta}$ و$\theta \in
\intoo{2\pi/3}{4\pi/3}$، يعطي السؤال 20
\[
PA = 2\,\abs{\sin\tfrac\theta2},
\quad
PB = 2\,\abs{\sin\bigl(\tfrac\theta2 - \tfrac\pi3\bigr)},
\quad
PC = 2\,\abs{\sin\bigl(\tfrac\theta2 - \tfrac{2\pi}3\bigr)} .
\]
وعلى القوس، $\frac\theta2 \in \intoo{\pi/3}{2\pi/3}$: عندئذ
$\sin\frac\theta2 > 0$؛ و$\frac\theta2 - \frac\pi3 \in
\intoo0{\pi/3}$، فيكون الجيب الثاني موجبًا؛ و$\frac\theta2 -
\frac{2\pi}3 \in \intoo{-\pi/3}0$، فيكون الثالث سالبًا و
$PC = 2\sin\bigl(\frac{2\pi}3 - \frac\theta2\bigr)$. وبتحويل المجموع إلى
جداء:
\[
\sin\Bigl(\frac\theta2 - \frac\pi3\Bigr) +
\sin\Bigl(\frac{2\pi}3 - \frac\theta2\Bigr)
= 2\,\sin\frac\pi6\,\cos\Bigl(\frac\theta2 - \frac\pi2\Bigr)
= \sin\frac\theta2 ,
\]
ومنه $PB + PC = 2\sin\frac\theta2 = PA$: وهي مبرهنة فان سخوتن.

\textbf{22.} عند $p = -1$: $PA = 2$ و $PB = PC = 1$، وطول كل ضلع من
المثلث المتساوي الأضلاع $\sqrt3$، فيصير طرفا
بطليموس $2\sqrt3$. وجبريًا، باستعمال $-1 - j^2 = j$ و
$1 + j = -j^2$:
\[
(a - b)(p - c) = (1 - j)\,(-1 - j^2) = (1 - j)j = j - j^2
= \iu\sqrt3 ,
\]
\[
(a - c)(b - p) = (1 - j^2)(j + 1) = 1 + j - j^2 - j^3
= j - j^2 = \iu\sqrt3 .
\]
والحدّان متساويان، فنسبتهما
$1 \in \R_{>0}$: أي حالة التساوي في السؤال 19، وهي تطابق $PA
\cdot BC = PB \cdot AC + PC \cdot AB$ تمامًا.

\textbf{23.} من أجل $(a, b, c) = (1, j, j^2)$: $n_a = \alpha j +
\beta j^2 = -\beta + \beta j^2$ (السؤال 14) $= \beta(j^2 - 1)$؛
وعدديًا $\beta(j^2 - 1) = \frac{(3 - \iu\sqrt3)}6 \cdot
\bigl(-\frac32 - \iu\frac{\sqrt3}2\bigr) = -1$. وبالمثل $n_b =
\alpha j^2 + \beta = -j$ و$n_c = \alpha + \beta j = -j^2$. إذن
مثلث نابليون الخارجي للمثلث $(1, j, j^2)$ هو $(-1, -j, -j^2)$: أي
المثلث الأصلي متناظرًا بالنسبة إلى مركز ثقله $0$ --- الحجم نفسه،
مدارًا بنصف دورة. فالمثلث المتساوي الأضلاع شكل صامد
لإنشاء نابليون، لا نهاية متقلصة.

\textbf{24.} (أ) الضرب في عدد أحاديّ \omterm{def:b1:complex:field}{المقياس} بوصفه دورانًا
هو ما بنى الرؤوس والمراكز (الأسئلة 2--3 و 11) وحوّل
المسافات على الدائرة إلى جيوب (السؤال 20). (ب) وبنية الزمرة
والشكل بوصفه ثابتًا هما ما برّر تسوية الدائرة
المحيطة إلى دائرة الوحدة وتسوية المثلث إلى $(1, j, j^2)$
في السؤال 21، عبر التساوق في السؤال 10. (ج) وتعميل
نصف الزاوية هو ما شغّل السؤال 20 وخطوةَ تحويل المجموع إلى
جداء التي أتمّت مبرهنة فان سخوتن.

\textbf{25.} يحوّل القاموس العبارات الهندسية إلى
متطابقات كثيرة الحدود يصير فيها كل فرض معادلة:
فما يربحه هو الميكنة --- إذ رُدّت مبرهنة نابليون إلى
$\alpha j + \beta = 0$، أي سطر واحد من الحساب في
$\Q(\iu\sqrt3)$؛ وما يكلّفه هو الوضوح الهندسي --- إذ إن
الحساب يشهد ولا يُظهر \emph{لماذا} تنغلق المراكز
في مثلث متساوي الأضلاع. والجواب المنصف أن
المتطابقة تشرح \emph{حتمية} المبرهنة (فهي تتحقق
تطابقيًا في $a, b, c$، فلا دخل لأيّ براعة في التشكيل)،
بينما يشرح الرسم \emph{مضمونها}. ويعود القاموس
نفسه في صورة مصفوفية عندما تصير الدورانات
المصفوفات المتعامدة $2 \times 2$ في \cref{ch:b1:matrices} و
\cref{ch:b1:euclid}، حيث يكون «الضرب في
$\eu^{\iu\theta}$» النموذجَ الأصلي للتقايس الخطي.
\end{solution}
